%-*-latex-*-
%-*-latex-*-
\newcommand\COURSE{ciss350}
\newcommand\ASSESSMENT{q10204}
\newcommand\ASSESSMENTTYPE{Quiz}
\newcommand\POINTS{\textwhite{xxx/xxx}}

\input{myquizpreamble}
\input{yliow}
\input{\COURSE}
%-*-latex-*-

%-*-latex-*-

%-*-latex-*-
\renewcommand\TITLE{\ASSESSMENTTYPE \ \ASSESSMENT}

\renewcommand\EMAIL{}
\renewcommand\AUTHOR{eagraham1@cougars.ccis.edu}

\textwidth=6in
\begin{document}
\topmatterthree

\LATEXHELPTHREEFIVEZERO

\TURNPAGE

%------------------------------------------------------------------------------
\newpage
\textsc{Compute best and worst runtime}

The goal is to compute the best and worst runtime for this algorithm:
\begin{Verbatim}[frame=single,fontsize=\small]
y = x[0]
z = x[1]
for i = 2, ..., n - 1:
    y = y + z
    if y + z + x[i] > 0:
        y = y + x[i - 2]
        z = x[i - 1]
    x[i - 2] = y
\end{Verbatim}
Clearly, the best case scenario is when the body of the if statement does not
execute
(for each iteration).
The worst case is when the body of the if statement executes executes
for each iteration.
First of all, here's the algorithm written in the simplest language
and with times attached.
Study it carefully, especially how to translate the if into
if-goto.
\begin{Verbatim}[frame=single,fontsize=\small]
          y = x[0]                  t1          
          z = x[1]                  t2
          i = 2                     t3
LOOP:     if i >= n:                t4
              goto ENDLOOP          t5
          y = y + z                 t6
          if y + z + x[i] <= 0:     t7
              goto ENDIF            t8
          y = y + x[i - 2]          t9
          z = x[i - 1]              t10
ENDIF:    x[i - 2] = y              t11
          i = i + 1                 t12
          goto LOOP                 t13
ENDLOOP:
\end{Verbatim}

\TURNPAGE
\newpage

%------------------------------------------------------------------------------
\nextq
For the best runtime,
include the number of times each statement is executed.
\\
\ANSWER
\begin{answercode}
          y = x[0]                  t1    1
          z = x[1]                  t2    1
          i = 2                     t3    1
LOOP:     if i >= n:                t4    n
              goto ENDLOOP          t5    1
          y = y + z                 t6    n - 2
          if y + z + x[i] <= y:     t7    n - 2
              goto ENDIF            t8    n - 2
          y = y + x[i - 2]          t9    0
          z = x[i - 1]              t10   0
ENDIF:    x[i - 2] = y              t11   n - 2
          i = i + 1                 t12   n - 2
          goto LOOP                 t13   n - 2
ENDLOOP:
\end{answercode}

%------------------------------------------------------------------------------
\nextq
Write down the best runtime function $T_b(n)$ as an expression involving
$n$ and constant times
$t_{1}$, $t_{2}$, ..., $t_{12}$.
Write $T_b(n)$ as a polynomial in descending power of $n$ -- i.e.,
the \textit{highest degree term appearing first}.
\\
\ANSWER
\begin{answerlong}
$T_b(n) = (t_{4} + t_{6} + t_{7} + t_{8} + t_{11} + t_{12} + t_{13})n + (t_{1} + t_{2} + t_{3} + t_{5} - 2t_{6} - 2t_{7} - 2t_{8} - 2t_{11} - 2t_{12} - 2t_{13})$
\end{answerlong}
(Correct the above)

%------------------------------------------------------------------------------
\nextq
Simplify the runtime function $T_b(n)$
by giving names $A$, $B$, $C$, $D$, ...
to the constants of the function from (a). (The $t_1, t_2, ...$ should
disappear.)
\\
\ANSWER
\begin{answerlong}
$T_b(n) = An + B$
\end{answerlong}
(Correct the above)

%------------------------------------------------------------------------------
\nextq
Fudge away the constants and write down the simplest $g(n)$ such
that the time in (b) is a big-O of your $g(n)$.
Your $g(n)$ should be either $n$ or $n^2$ or $n^3$ or $n^{102.25}$ ...
\\
\ANSWER
\begin{answerlong}
$T_b(n) = O(n)$
\end{answerlong}
(Correct the above)

%------------------------------------------------------------------------------
\nextq
Second, write down the worst runtime of the above algorithm.
(You need not go through all the computations. You just need to state
the worst runtime.)
\\
\ANSWER
\begin{answerlong}
$T_w(n) = O(n)$
\end{answerlong}

%------------------------------------------------------------------------------
%-*-latex-*-
\newpage
\input{instructions}
\end{document}

